\chapter{Quaternionic Preliminaries}
\label{app:quaternions}

Chapters~\ref{ch:quaternionic} and~\ref{ch:quaternionic-mub} used
quaternionic matrix algebra without formally collecting its conventions in
one place. This appendix does so, stating precisely the facts those
chapters relied on — including the two that are \emph{not} true by naive
analogy with the complex case, which is exactly where the earlier chapters
found real difficulty.

\section{The Quaternion Algebra}

A quaternion is $q = a+bi+cj+dk$ with $a,b,c,d\in\R$ and $i^2=j^2=k^2=
ijk=-1$. Quaternion multiplication is associative but not commutative
($ij=k \ne -k = ji$). The conjugate is $\bar q = a-bi-cj-dk$, and the norm
is $|q|^2 = q\bar q = a^2+b^2+c^2+d^2 \in \R_{\ge0}$. Every nonzero
quaternion has a multiplicative inverse $q^{-1}=\bar q/|q|^2$ — the norm's
multiplicativity, $|pq|=|p||q|$, is what makes $\Hset$ a genuine division
algebra, not merely a ring. The group
of unit quaternions ($|q|=1$) is denoted $\mathrm{Sp}(1)$; it is a
three-real-dimensional Lie group, isomorphic to $SU(2)$, in contrast to
the complex case's one-dimensional $U(1)$.

\begin{fact}[Frobenius' theorem]
$\R$, $\C$, and $\Hset$ are the only finite-dimensional associative
division algebras over $\R$. Relaxing associativity admits exactly one
further case, the octonions $\mathbb O$ (not used anywhere in this book).
\end{fact}

\begin{fact}[Scalar--vector decomposition]
\label{fact:scalar-vector}
Writing $q=(r,\mathbf v)$ with $r\in\R$ and $\mathbf v\in\R^3$ (identifying
$bi+cj+dk$ with $(b,c,d)$), quaternion multiplication becomes
\begin{equation}
(r_1,\mathbf v_1)(r_2,\mathbf v_2) = \big(r_1r_2 - \mathbf v_1\cdot\mathbf
v_2,\ \ r_1\mathbf v_2+r_2\mathbf v_1+\mathbf v_1\times\mathbf v_2\big),
\end{equation}
verified directly by expanding both sides in $i,j,k$. The familiar dot and
cross products of $\R^3$ are not independent constructions; they are the
symmetric and antisymmetric parts of a single quaternion product.
\end{fact}

\section{Unit Quaternions and Rotations}

\begin{fact}[The rotation formula]
\label{fact:rotation-formula}
For a unit vector $\mathbf n\in\R^3$ and angle $\theta$, let $q =
\cos\theta + \mathbf n\sin\theta \in \mathrm{Sp}(1)$ (using
Fact~\ref{fact:scalar-vector}'s identification). For a pure-vector
quaternion $v=(0,\mathbf v)$, $v \mapsto qvq^{-1}$ rotates $\mathbf v$
about axis $\mathbf n$ by angle $2\theta$.
\end{fact}

This was checked computationally (not merely cited): for $\mathbf n$ the
$z$-axis and $\theta=\pi/6$, $qvq^{-1}$ applied to the unit vector along
$x$ reproduces the standard $60^\circ$ rotation about $z$ exactly.

\begin{fact}[$\mathrm{Sp}(1)$ double-covers $SO(3)$]
The map $q \mapsto (\mathbf v \mapsto qvq^{-1})$ sends $\mathrm{Sp}(1) \cong
SU(2) \cong S^3$ onto $SO(3)$, two-to-one: $q$ and $-q$ act identically on
every vector, since $(-q)v(-q)^{-1} = qvq^{-1}$. This is the same
$\mathrm{Sp}(1)$ that appeared as the local phase group in
Chapter~\ref{ch:quaternionic-mub}, and the same double-cover structure
underlies the $720^\circ$ return of spin-$\tfrac12$ systems in ordinary
quantum mechanics — a fact this book does not derive or otherwise depend
on, recorded here only because the same group appeared for an unrelated
reason (MUB counting) and the coincidence is worth naming.
\end{fact}

\section{The Clifford Algebra Correspondence}

\begin{fact}
\label{fact:clifford}
$\Hset$ is isomorphic to the even subalgebra of the Clifford algebra of
$\R^3$: $\Hset \cong Cl^+_{3,0}(\R)$. Under this isomorphism the basis
elements $i,j,k$ correspond to the three coordinate bivectors $e_2e_3,
e_3e_1, e_1e_2$ — oriented coordinate planes, not an extra spatial
direction. This is standard material in geometric (Clifford) algebra,
not a claim original to this book.
\end{fact}

\begin{remark}
Fact~\ref{fact:clifford} reframes what "imaginary" means for quaternions:
$i,j,k$ generate rotations in the $yz$, $zx$, $xy$ planes respectively,
and quaternion multiplication is fundamentally about composing such
rotations rather than an extension of ordinary arithmetic by a mysterious
extra unit. This is consistent with, and gives a coordinate-free gloss on,
Chapter~\ref{ch:amplitwist}'s identification of the complex unit with the
generator of a plane rotation ($J = R(\pi/2)$): the quaternionic case is
the same idea with three independent planes instead of one.
\end{remark}

\section{Quaternionic Matrices}

A quaternionic $n\times n$ matrix $M$ has entries in $\Hset$; matrix
multiplication is defined exactly as over a commutative ring, and remains
associative (inherited from associativity of quaternion multiplication)
though the entries do not commute. The conjugate transpose is
$M^\dagger_{ij} := \overline{M_{ji}}$; $M$ is Hermitian if $M^\dagger=M$,
which forces every diagonal entry to be real (a self-conjugate quaternion
is real). $M$ is unitary if $M^\dagger M = I$.

\begin{fact}[Associativity suffices for graph propagation]
Every argument in Chapter~\ref{ch:context-graphs}--
\ref{ch:minimal-witnesses} building a global frame from local transitions
(tree propagation, cycle closure) uses only associativity of matrix
multiplication together with the existence and well-definedness of
$\pi(L)_{ij} := |L_{ij}|^2_\Hset = L_{ij}\overline{L_{ij}}$. Neither
requires commutativity, so these arguments transport to $\Hset$ unchanged
(Chapter~\ref{ch:quaternionic}, Proposition~\ref{prop:h-generalizes}).
\end{fact}

\section{Two Traps}

\begin{fact}[Trace is not cyclic]
\label{fact:trace-noncyclic}
For general quaternionic matrices $A,B$, $\operatorname{Tr}(AB) \ne
\operatorname{Tr}(BA)$ — verified directly by example in
Chapter~\ref{ch:quaternionic-mub}. This is the genuine obstruction to a
naive generalization of complex trace-based arguments, not the
$\mathrm{Sp}(1)$-versus-$U(1)$ phase freedom originally suspected to be the
difficulty (Chapter~\ref{ch:quaternionic}'s closing remark).
\end{fact}

\begin{fact}[The real part is cyclic]
\label{fact:trace-real-cyclic}
Despite Fact~\ref{fact:trace-noncyclic}, $\operatorname{Re}[\operatorname{Tr}
(AB)] = \operatorname{Re}[\operatorname{Tr}(BA)]$ always. This makes
$\langle A,B\rangle_\Hset := \operatorname{Re}[\operatorname{Tr}(A^\dagger
B)]$ the correct symmetric real inner product on the space of
quaternionic matrices, and it is this inner product — not the raw
quaternionic trace — that the orthogonality argument underlying
Theorem~\ref{thm:h-mub-bound} depends on
(Proposition~\ref{prop:h-real-part}).
\end{fact}

\begin{fact}[Projectors are $\mathrm{Sp}(1)$-phase-invariant]
For a unit quaternion $q$ and a normalized $|e\rangle \in \Hset^d$,
$(|e\rangle q)(|e\rangle q)^\dagger = |e\rangle\langle e|$: the rank-one
projector built from a state vector does not depend on which unit
quaternion represents its ray, exactly as the complex case's $U(1)$ phase
drops out of $|\psi\rangle\langle\psi|$. Consequently the extra two real
dimensions of $\mathrm{Sp}(1)$ freedom (compared to $U(1)$) contribute
nothing to the space spanned by one orthonormal basis's traceless
projectors, which remains $(d-1)$-dimensional exactly as in the complex
case (Proposition~\ref{prop:h-di-1}). This was the concern flagged when
Chapter~\ref{ch:quaternionic} was first written; it resolves harmlessly,
and the genuine difficulty was Fact~\ref{fact:trace-noncyclic} instead.
\end{fact}

\section{Summary Table}

\begin{center}
\begin{tabular}{@{}lll@{}}
\toprule
& Complex ($\C$) & Quaternionic ($\Hset$) \\
\midrule
Local phase group & $U(1)$, 1 real dim. & $\mathrm{Sp}(1)$, 3 real dim. \\
Matrix multiplication & associative, commutative & associative, non-commutative \\
Trace cyclicity & $\operatorname{Tr}(AB)=\operatorname{Tr}(BA)$ & fails; $\operatorname{Re}[\operatorname{Tr}(AB)]=\operatorname{Re}[\operatorname{Tr}(BA)]$ only \\
$\beta_{\mathbb F}(d)$ & $d^2$ & $d(2d-1)$ \\
MUB bound & $d+1$ & $2d+1$ (Theorem~\ref{thm:h-mub-bound}) \\
\bottomrule
\end{tabular}
\end{center}
